\chapter{Correttezza}

Il test di correttezza è stato implementato in \texttt{C++} all'interno del file \texttt{correctness.cpp}.

Sono state considerate quattro diverse tipologie di dati:
\texttt{integer}, \texttt{random strings}, \texttt{real words} e \texttt{hybrid}, 
dove la chiave è di tipo \texttt{real word} e il valore è di tipo \texttt{integer}.

Nel file è stata implementata una sezione di test per ognuna delle quattro tipologie di dato, 
con l'obiettivo di verificare la correttezza delle operazioni fondamentali e di collezionare 
eventuali errori riscontrati durante l'esecuzione.

\begin{lstlisting}[language=C++, caption=Integer test]
// Integer
{
    int errors = 0;
    for(int x=0; x<20; x++) {
        for(int i=0; i<6; i++) {
            int current_error = intTest(std::pow(10, i));
            std::cout << "Partial Operation Failed: " << current_error << std::endl;
            std::cout << "-----------------------------" <<std::endl;
            errors += current_error;
        }
    }
    arr_error[0] += errors;
    std::cout << "-----------------------------" <<std::endl;
    std::cout << "-----------------------------" <<std::endl;
    std::cout << "INTEGER ERRORS: " << errors << std::endl;
    std::cout << "-----------------------------" <<std::endl;
    std::cout << "-----------------------------" <<std::endl;
}
\end{lstlisting}

Ogni tipologia di dato utilizza una funzione di test dedicata, che si occupa di verificare 
il comportamento della struttura dati su diverse configurazioni di hash table. 
Per ciascuna configurazione viene creata una nuova istanza della classe \texttt{HashTable}, 
parametrizzata con le specifiche modalità di gestione delle collisioni, funzione di hash e strategia di ridimensionamento.

Il test viene inoltre ripetuto più volte, variando il numero di elementi utilizzati, 
al fine di ridurre l’impatto della casualità e ottenere risultati più affidabili.

\begin{lstlisting}[language=C++, caption=Integer test function]
/*
* HashTable:
*   OVERFLOW::OPEN_ADDRESSING_LINEAR_PROBING
*   HASH::MODE_1
*   RESIZE::DOUBLE
* 
* Key::INTEGER
* Value::INTEGER
*/
{
    std::cout << "-----------------------------" <<std::endl;
    std::cout << "HashTable:" << std::endl;
    std::cout << "  OVERFLOW::OPEN_ADDRESSING_LINEAR_PROBING" << std::endl;
    std::cout << "  HASH::MODE_1" << std::endl;
    std::cout << "  RESIZE::DOUBLE" << std::endl;
    std::cout << "  Key::INTEGER" << std::endl;
    std::cout << "  Value::INTEGER" << std::endl;
    std::cout << "-----------------------------" <<std::endl;

    HashTable<int, int> HashTable(OVERFLOW::OPEN_ADDRESSING_LINEAR_PROBING, HASH::MODE_1, RESIZE::DOUBLE);
    errors += intTestRoutine(HashTable, elements);
}
\end{lstlisting}

Per ogni configurazione viene eseguita una routine specifica, differenziata in base alla tipologia di dato, 
che verifica la correttezza delle operazioni fondamentali della hash table.

In particolare:
\begin{itemize}
    \item \textbf{Insert}: viene verificato che ogni elemento inserito sia effettivamente memorizzato nella struttura e accessibile 
    tramite operazione di ricerca (\texttt{find}).
    \item \textbf{Update}: viene controllato che i valori associati alle chiavi siano correttamente aggiornati.
    \item \textbf{Remove}: viene verificato che gli elementi rimossi non siano più accessibili e che l’operazione non comprometta l’accesso agli altri 
    dati presenti nella struttura.
\end{itemize}

\begin{lstlisting}[language=C++, caption=Integer test routine for insert operation]
const int ELEMENTS = elements;
int errors = 0;
// Generatore di numeri casuali
std::vector<std::pair<int, int>> numbers;
numbers.reserve(ELEMENTS);

std::random_device rd;
std::mt19937 gen(rd());
std::uniform_int_distribution<> dist(1, ELEMENTS*10);

int value = 0;
for (int i = 0; i < ELEMENTS; ++i) {
    value = dist(gen);
    numbers.emplace_back(value, value);
}

// Inserts
{
    std::cout << "-----------------------------" <<std::endl;
    std::cout << "INSERT" <<std::endl;
    std::cout << "-----------------------------" <<std::endl;
    std::vector<std::pair<int, int>> failed;
    failed = insert(HashTable, numbers);
    errors += failed.size();
    // Exist
    failed = exist(HashTable, numbers);
    errors += failed.size();
    // Get
    failed = get(HashTable, numbers);
    errors += failed.size();
}
\end{lstlisting}

Durante l’esecuzione dei test, il numero di errori riscontrati viene accumulato e riportato per ciascuna tipologia di dato e configurazione testata.

L’esecuzione complessiva dei test non ha evidenziato errori, confermando la correttezza dell’implementazione delle operazioni fondamentali 
della hash table nelle diverse configurazioni considerate.

\begin{lstlisting}[language=C++, caption=Correctness test result]
-----------------------------
-----------------------------
TOTAL ERRORS: 0
-----------------------------
INTEGER ERRORS: 0
STRING ERRORS: 0
REAL WORD ERRORS: 0
HYBRID ERRORS: 0
-----------------------------
-----------------------------
\end{lstlisting}